\documentclass[11pt]{article}
\usepackage[paperwidth=210mm,paperheight=297mm,textwidth=182mm,top=20mm,bottom=22mm]{geometry}
\usepackage{amsmath,amssymb,newtxmath}
\usepackage[no-math]{fontspec}
\setmainfont{Times New Roman}
\usepackage{xeCJK}
\setCJKmainfont[AutoFakeBold=2]{SimSun}
\usepackage{graphicx,booktabs,array,hyperref,placeins}
\hypersetup{hidelinks,unicode=true}
\raggedbottom
\setlength{\emergencystretch}{2em}
\renewcommand{\arraystretch}{1.15}
\usepackage{seqsplit}
\begin{document}
\title{Litz 线分层建模与数值验证：完整研究报告}
\author{缪雨轩\\淮南师范学院}
\date{}
\maketitle
\FloatBarrier
\section*{摘要}
\hypertarget{abstract}{}
\pdfbookmark[0]{摘要}{bookmark-abstract}
Litz 线的高频损耗取决于单股趋肤、局部邻近场、股间复电流分配、换位路径及端部约束。本文建立可复现的分层数值验证框架，区分各层实际保留的物理内容。沿途绝缘、端部并联的股线采用共同端电压和总复电流约束，并在指定线性场路模型内给出均流判据。有限长度纵向 PEEC 则独立检查轴向长度、互项与自项积分及横向离散的影响。

在 200 kHz、20 A RMS、20 ℃ 铜导体工况下，单圆线解析与独立二维校准建立了内场参考，640 股二维共端模型具有受控网格、回流和外域检查。5024 自由度纵向模型末档 Rac 为 20.012738 mΩ/轴向m，与本档等面积 Bessel 参照相差 0.044582\%。固定长方体外场 RT0 算例的末两次损耗变化为 0.655090\% 和 0.362916\%；它检验外场算子，不能作为 Litz 输运电阻。三维 A–φ、Ansys 新接口及真实性能排名的未通过边界继续保留。结果支持限定模型内的数值结论，不证明实际多级或跨层 Litz 的性能优劣。


关键词：Litz 线；分层验证；趋肤与邻近效应；共端场路；有限长度 PEEC；可复现数值研究


\FloatBarrier
\section*{01 引言与研究问题}
\hypertarget{chapter-01}{}
\pdfbookmark[0]{01 引言与研究问题}{bookmark-chapter-01}
\subsection*{引言}
\hypertarget{01.introduction.heading}{}
\pdfbookmark[1]{引言}{bookmark-01.introduction.heading}
将总铜截面分成许多细股，并不能单独保证高频损耗下降。单股直径决定局部趋肤尺度，邻近导体与回流决定外部磁场，沿程位置和共同端接则决定各股复电流。仅比较几何外观、名义铜面积或平均电流，会把不同物理机制混合在同一个指标中。可靠的结构比较应首先确定哪些场量和约束被模型保留，再检查相应的数值不确定性。


现有方法提供了不同尺度的处理途径。离散换位与二维 FEM 相结合的 DTM [1] 将位置置换与截面场解联系起来；非理想复合换位 iDTM [2] 进一步关注层级及子束换位。iSEC [3] 在其指定绕组与分层假设下构造降阶电路；FEM–PEEC 耦合 [4] 和准二维多导体方法 [5] 则采用不同的宏观场、端场与电流约束。连续体模型 [6]、复磁导率均匀化 [7] 降低了大规模绕组的计算成本，但其平均场不能直接替代逐股复电流或真实端区。


本文的重点不是提出已经得到验证的最优换位结构，而是建立结构研究所需的可比较量与验证链。第一项工作是区分解析、二维、纵向和完整体素模型的验证关系；第二项是从共端增广系统推导均流条件，明确其模型范围；第三项是对有限长度纵向 PEEC 进行单因素敏感性检查，并把 CPU/GPU 算法对照与物理验证分别解释。传统多级 Litz 保持强基线；跨层几何作为待验证候选，不预设优越性。


\subsection*{研究问题与贡献}
\hypertarget{01.scope.heading}{}
\pdfbookmark[1]{研究问题与贡献}{bookmark-01.scope.heading}
本研究围绕三个问题组织验证：局部趋肤基准能否由独立离散方法复现；共同端电压下的复电流分配由哪些代数条件决定；有限长度纵向核的电阻结果对长度和积分设置有多大敏感性。所给出的贡献是统一物理量、可审查的验证关系和模型内的定量证据，供真实换位结构研究选择合适的下一层模型。


\begin{figure}[!htbp]
\centering
\setcounter{figure}{0}
\refstepcounter{figure}\label{fig:model_hierarchy}
\hypertarget{fig-model_hierarchy}{}
\includegraphics[width=88.9mm]{../../../panels/zh-CN/model_hierarchy.pdf}
\par\vspace{4pt}{\small\centering 图 1. 分层模型与可比较观测量。\par}
\vspace{3pt}{\small\raggedright 输运模型按物理对象与观测量分组：圆线内场用于解析与二维校准，多股二维场路求解共同端口下的复电流，纵向 PEEC 检验有限长度效应。连线标明内部场校准或等面积电阻参照。RT0 外场损耗与三维 A–φ 场诊断各自构成独立研究范围。\par}
\end{figure}
\FloatBarrier
\section*{02 物理约定与分层建模框架}
\hypertarget{chapter-02}{}
\pdfbookmark[0]{02 物理约定与分层建模框架}{bookmark-chapter-02}
\subsection*{工况、物理量与公平比较}
\hypertarget{02.conventions.heading}{}
\pdfbookmark[1]{工况、物理量与公平比较}{bookmark-02.conventions.heading}
计算采用 SI 单位、exp(jωt) 和 RMS 相量。材料为 20 ℃ 铜，σ=5.8×10⁷ S/m、μr=1；主输运工况 f=200 kHz、总电流 20 A RMS，无独立外加场。股线沿途绝缘、端部并联。所有“每米”指轴向长度，实际股线弧长另行报告。有限段端电压 V 与二维模型的单位轴向长度端场具有不同单位，比较时必须保留长度归一化。


原题 Q1 的 D=2 mm 圆线保留为解析校准，其铜面积为 3.141593 mm²、Jdc=6.366198 A/mm²，超过 4 A/mm² 约束。它不是满足最终设计条件的导体。Q2 的名义铜面积至少 5 mm²；640×0.100 mm 的理想圆面积为 5.026548 mm²，而原二维 64 边形离散面积约 5.018478 mm²。采用实际面积的等铜参照，才能避免把几何面积误差解释成纯电磁改善。


不同股长的直流电阻由并联电导相加得到 \href{report_zh-CN_single.pdf#E01}{(1)}。同时报告铜体积、弧长增长、绝对 Rac、Rac/Rdc、损耗和逐股复电流，而不是只给比值。RMS 焦耳损耗采用 \href{report_zh-CN_single.pdf#E02}{(2)}，没有峰值相量的 1/2 因子。存在独立外场时，吸收功率含外场做功；RT0 外场损耗不能无条件除以 20² 命名为内禀输运电阻。几何等铜也不意味着单位轴向长度的铜体积相同。


\subsection*{磁准静态物理约定}
\hypertarget{02.maxwell.heading}{}
\pdfbookmark[1]{磁准静态物理约定}{bookmark-02.maxwell.heading}
在磁准静态近似下，位移电流相对于导电电流忽略，电荷连续性给出 \(\nabla\)·J=0。采用 exp(jωt) 后，法拉第定律为 \(\nabla\)×E=−jωB，欧姆关系为 J=σE。通过材料、激励、回流和边界共同定义一个边值问题，再将可比较的观测量映射到下一层离散问题，构成本文的验证顺序。


\subsection*{分层模型与验证方法}
\hypertarget{02.hierarchy.heading}{}
\pdfbookmark[1]{分层模型与验证方法}{bookmark-02.hierarchy.heading}
模型层级按未知量与可比较观测量划分。解析圆线给出长直均质导体的内部阻抗和径向电流；二维扩散验证相同内场问题。二维共端多股模型进一步保留股内分布、股间复电流及明确回流，但省略纵向换位与有限端区。纵向 PEEC 保留有限长度部分电感核和共端约束，却不包含完整横向环流。三分量 RT0 外场算例检验三维算子、连续性和外场吸收，仍与曲线 Litz 的输运问题不同。


验证同时包括求解残差、约束、独立功率与参数敏感性。真实线性残差和公共端/总电流约束沿用 10⁻⁸ 门槛，独立功率平衡为 0.1\%，相关网格、分段、边界和长度变化为 1\%。解析参考及解析电流/功率积分有更严格的专门门槛；RT0 GMRES 目标保持 10⁻¹¹，积分敏感性为 0.1\%；CPU/GPU 同矩阵与复电流对照为 10⁻¹⁰。这些门槛属于本研究的检查定义，不是全部物理问题已经具有同样精度的宣言。


证据可在共同假设下交叉比较，不能沿箭头自动升级结论。小线性残差说明已装配系统被正确求解；它不排除遗漏物理或错误边界。功率平衡约束离散方程与后处理的一致性，但保守的错误模型也可能平衡。两档结果接近限制特定扰动下的变化，不能给出无条件误差上界或统计置信区间。归档因此把执行状态、数值检查状态和物理适用范围分别保存。


\hypertarget{E01}{}
\begin{equation}
\begin{gathered}
R_{\mathrm{dc}}=\left(\sum_{k=1}^{N}\frac{\sigma A_k}{\ell_k}\right)^{-1} \\
V_{\mathrm{Cu}}=\sum_{k=1}^{N}A_k\ell_k
\end{gathered}
\tag{1}\label{eq:E01}
\end{equation}
\hypertarget{E02}{}
\begin{equation}
\begin{gathered}
P_{\mathrm{Cu}}=\int_{\Omega_{\mathrm{Cu}}}|\mathbf J|^2/\sigma\,\mathrm d\Omega \\
R_{\mathrm{ac}}=P_{\mathrm{Cu}}/|I_t|^2 \\
r_{\mathrm{ac}}=R_{\mathrm{ac}}/L_z
\end{gathered}
\tag{2}\label{eq:E02}
\end{equation}
\FloatBarrier
\section*{03 单圆线解析与二维校准}
\hypertarget{chapter-03}{}
\pdfbookmark[0]{03 单圆线解析与二维校准}{bookmark-chapter-03}
\subsection*{径向方程与归一化}
\hypertarget{03.round_derivation.heading}{}
\pdfbookmark[1]{径向方程与归一化}{bookmark-03.round_derivation.heading}
圆线内部的轴向电流满足 (1/r) d(r dJz/dr)/dr=κ²Jz。轴心的有限性选择 I\(_0\)(κr)；使用 \(\int\)\(_0\)\(^a\) rI\(_0\)(κr)dr=aI\(_1\)(κa)/κ 和总复电流约束即可确定归一化常数。表面轴向电场与总电流之比给出内部阻抗，其实部通过独立焦耳积分复核。该推导把解析形状、归一化和功率三个检查连接起来。


\subsection*{单圆线解析与独立二维校准}
\hypertarget{03.analytic.heading}{}
\pdfbookmark[1]{单圆线解析与独立二维校准}{bookmark-03.analytic.heading}
长直圆线内部解采用修正 Bessel 函数 \href{report_zh-CN_single.pdf#E03}{(3)}。200 kHz 下的趋肤深度 δ=0.147771655 mm，相对 1 mm 半径已形成明显表层电流集中。解析值为 Rdc=5.488101486 mΩ/m、Rac=20.016203628 mΩ/m、Rac/Rdc=3.647199980、损耗 8.006481451 W/m。这里的内部阻抗不含由回流确定的外部电感。径向相位保留原复相量，并通过相邻半径的 360° 连续提升显示；显示解缠不修改数据。


独立二维圆盘扩散使用三档 353、1389、5573 节点，并按总复电流归一化。最细 Rac 解析相对误差为 1.79233×10⁻⁶，功率积分差为 3.06800×10⁻⁶。结果支持该孤立内场问题的离散与归一化。二维场与有限长度三维端区不具有相同边界，不能把这个误差直接用于后者。解析截面图由径向解析幅值线性插值并轴对称映射。


\hypertarget{E03}{}
\begin{equation}
\begin{gathered}
\delta=(\pi f\mu\sigma)^{-1/2} \\
\kappa=(1+j)/\delta \\
J_z(r)=\frac{I_t\kappa I_0(\kappa r)}{2\pi a I_1(\kappa a)} \\
z_{\mathrm{int}}=\frac{\kappa I_0(\kappa a)}{2\pi a\sigma I_1(\kappa a)}
\end{gathered}
\tag{3}\label{eq:E03}
\end{equation}
\begin{figure}[!htbp]
\centering
\setcounter{figure}{1}
\refstepcounter{figure}\label{fig:q1_profile}
\hypertarget{fig-q1_profile}{}
\includegraphics[width=182mm]{../../../composites/zh-CN/q1_profile.pdf}
\par\vspace{4pt}{\small\centering 图 2. 单圆线径向电流幅值与相位。\par}
\vspace{3pt}{\small\raggedright 铜导体直径 2 mm，200 kHz、20 A RMS。 (a) RMS 幅值；(b) 连续显示的相位。 相位显示采用径向连续提升，幅值保留 RMS 归一化。径向样本来自保存的 Bessel 解析数据。\par}
\end{figure}
\begin{figure}[!htbp]
\centering
\setcounter{figure}{2}
\refstepcounter{figure}\label{fig:q1_cross_section}
\hypertarget{fig-q1_cross_section}{}
\includegraphics[width=88.9mm]{../../../panels/zh-CN/q1_cross_section.pdf}
\par\vspace{4pt}{\small\centering 图 3. 单圆线解析截面幅值。\par}
\vspace{3pt}{\small\raggedright 铜导体直径 2 mm，200 kHz、20 A RMS。 色标显示轴向 RMS 电流密度。 径向样本经线性插值后作轴对称映射。坐标以毫米显示，电流密度以 A/mm² 显示。\par}
\end{figure}
\begin{figure}[!htbp]
\centering
\setcounter{figure}{3}
\refstepcounter{figure}\label{fig:q1_fem_error}
\hypertarget{fig-q1_fem_error}{}
\includegraphics[width=88.9mm]{../../../panels/zh-CN/q1_fem_error.pdf}
\par\vspace{4pt}{\small\centering 图 4. 单圆线二维校准误差。\par}
\vspace{3pt}{\small\raggedright 铜导体直径 2 mm，200 kHz、20 A RMS。 三档网格与独立功率积分对应同一内场问题。 电阻参照为相同圆线的内部 Bessel 解。功率使用独立材料域积分计算。\par}
\end{figure}
\FloatBarrier
\section*{04 多股共端场路与均流判据}
\hypertarget{chapter-04}{}
\pdfbookmark[0]{04 多股共端场路与均流判据}{bookmark-chapter-04}
\subsection*{共端场路与复电流均流判据}
\hypertarget{04.common.heading}{}
\pdfbookmark[1]{共端场路与复电流均流判据}{bookmark-04.common.heading}
端部并联并不要求各股预先具有相同电流。令 Z 为指定线性模型的耦合阻抗矩阵，e 为独立外场的感应项，i 为股线复电流向量，则公共端电压和总电流形成增广系统 \href{report_zh-CN_single.pdf#E04}{(4)}。实际实现求解增广系统，不依赖显式矩阵求逆。约束允许幅值和相位由场路共同决定；仅比较 |ik| 而忽略相位，会漏掉复电流和总电流之间的重要关系。


在增广系统解唯一的前提下，把 i=(It/N)1 代入端口方程，均流成立当且仅当 [(It/N)Z1+e] 落在 1 的张成空间。用投影 Q 表示即得到 \href{report_zh-CN_single.pdf#E05}{(5)}。无独立外场时，这成为阻抗行和相同的条件；外场存在时，感应项可能保持或破坏该条件。这个代数判据适用于所指定的 Z、e 和端接，不是任意连续三维换位结构的均流定理。


二维共端 FEM 将公共轴向端场与 Az 自由度一起求解。保存的 u 与端场具有 V/m 单位，Jz=σ(v−u)；回流使用自己的端场，空气不赋导体电流。图件从原节点和材料三角形恢复 Jz，按原毫米坐标缩放执行积分。独立线性三角形质量积分与各股复电流、束内损耗、回流损耗和铜面积核对，因此颜色图的来源可审查。保存的 u 不能直接当作电流密度绘制。


\subsection*{判据推导}
\hypertarget{04.sharing_derivation.heading}{}
\pdfbookmark[1]{判据推导}{bookmark-04.sharing_derivation.heading}
均流判据的充分性来自端口解的唯一性。当投影条件成立时，向量 (It/N)Z1+e 的各分量相同，因此存在共同电压使等流向量满足全部端口方程及总电流约束。必要性则由已存在的等流解直接代入得到。数值计算以归一化复电流偏差评估均流程度，同时报告幅值和相位，以保持复数求和的信息。


\subsection*{640 股二维结果及其边界}
\hypertarget{04.coax_results.heading}{}
\pdfbookmark[1]{640 股二维结果及其边界}{bookmark-04.coax_results.heading}
640 股直束的冻结三档网格、回流与外域检查共有 23 项通过。最细网格增量的 Rac 变化为 0.011902\%，回流距离变化为 0.000495\%，外域变化为 0.000270\%，均低于本路线 1\% 门槛。它们说明指定二维共端模型在已检查扰动范围内稳定，不能代表真实绕制束的有限端区已经收敛。束内 Rac 与回流损耗分别报告，避免把回流中的热损耗算成导体自身改善。


逐股幅值与相位展示了共同端压下的电流分配，虚线 It/N 仅为参照，不是求解约束。节点电流积分与公共端条件的一致性，是均流讨论的必要前提。对这种二维直束，把单股直径减小并不等价于沿程位置已经换位。它的主要用途是验证多股场路与受控边界，作为后续换位模型的截面证据；最终结构排序仍需要保留沿程耦合、实际端接与相关数值不确定性的模型或实验。


\subsection*{束内与回流能量}
\hypertarget{04.interpretation.heading}{}
\pdfbookmark[1]{束内与回流能量}{bookmark-04.interpretation.heading}
640 股截面结果说明了细股自身趋肤尺度与束内集体电流分配之间的区别。外层与内层股线具有不同的磁耦合环境，共同端电压使这种差异表现为复电流的幅值和相位变化。此时直流面积给出铜利用率的基准，束内焦耳积分给出输运损耗，而回流的单独积分解释完整场路的能量分配。


640 股束内 Rac=15.793003197 mΩ/轴向m，束内损耗 6.317201279 W/轴向m，回流损耗 2.196205516 W/轴向m，实际铜截面 5.018478 mm²。共同端轴向电场为 0.411856245+j(4.064969103) V/m。归一化复电流 RMS 偏差为 1.854214，各股电流的复数和恢复 20 A RMS 总输运。


\hypertarget{E04}{}\begin{equation}\begin{bmatrix} Z&-\mathbf1\\ \mathbf1^T&0\end{bmatrix}\begin{bmatrix}\mathbf i\\V\end{bmatrix}=\begin{bmatrix}-\mathbf e\\I_t\end{bmatrix}\tag{4}\label{eq:E04}\end{equation}
\hypertarget{E05}{}
\begin{equation}
\begin{gathered}
Q=I_N-\mathbf1\mathbf1^T/N \\
Q\left[(I_t/N)Z\mathbf1+\mathbf e\right]=\mathbf0
\end{gathered}
\tag{5}\label{eq:E05}
\end{equation}
\begin{figure}[!htbp]
\centering
\setcounter{figure}{4}
\refstepcounter{figure}\label{fig:common_terminal}
\hypertarget{fig-common_terminal}{}
\includegraphics[width=88.9mm]{../../../panels/zh-CN/common_terminal.pdf}
\par\vspace{4pt}{\small\centering 图 5. 多股共同端口与回流。\par}
\vspace{3pt}{\small\raggedright 各股沿途绝缘，在两端并联，逐股复电流由共同端电压及总电流约束求解。耦合矩阵 Z 与激励项 e 包含指定模型中的阻抗和场作用。受控回流承担反向总电流，并具有独立的端场。图中给出已知输入及待求复电流和电压。\par}
\end{figure}
\begin{figure}[!htbp]
\centering
\setcounter{figure}{5}
\refstepcounter{figure}\label{fig:coax640_currents}
\hypertarget{fig-coax640_currents}{}
\includegraphics[width=182mm]{../../../composites/zh-CN/coax640_currents.pdf}
\par\vspace{4pt}{\small\centering 图 6. 640 股共端复电流分配。\par}
\vspace{3pt}{\small\raggedright (a) 幅值；(b) 相位，按物理股线身份排列。 虚线表示 It/N 的比较参照。幅值和相位使用相同的物理股线编号。\par}
\end{figure}
\begin{figure}[!htbp]
\centering
\setcounter{figure}{6}
\refstepcounter{figure}\label{fig:coax640_fields}
\hypertarget{fig-coax640_fields}{}
\includegraphics[width=182mm]{../../../composites/zh-CN/coax640_fields.pdf}
\par\vspace{4pt}{\small\centering 图 7. 640 股二维截面电流密度。\par}
\vspace{3pt}{\small\raggedright (a) RMS 幅值；(b) 相位，材料标记识别各股。 电流按 Jz=σ(v−u) 从原节点场与端场恢复。独立功率使用线性三角形质量积分，显示采用单元中心采样。\par}
\end{figure}
\begin{figure}[!htbp]
\centering
\setcounter{figure}{7}
\refstepcounter{figure}\label{fig:coax640_mesh}
\hypertarget{fig-coax640_mesh}{}
\includegraphics[width=88.9mm]{../../../panels/zh-CN/coax640_mesh.pdf}
\par\vspace{4pt}{\small\centering 图 8. 640 股网格与边界敏感性。\par}
\vspace{3pt}{\small\raggedright 各试验按网格、回流和外域分别改变一个因素。 变化以相邻指定工况的电阻差归一化。网格、回流距离与外域大小各使用独立序列。\par}
\end{figure}
\FloatBarrier
\section*{05 换位模型、真实几何与等铜比较}
\hypertarget{chapter-05}{}
\pdfbookmark[0]{05 换位模型、真实几何与等铜比较}{bookmark-chapter-05}
\subsection*{几何、实际弧长与离散换位}
\hypertarget{05.geometry.heading}{}
\pdfbookmark[1]{几何、实际弧长与离散换位}{bookmark-05.geometry.heading}
保存的直束、单级、传统多级和跨层中心线用于检查股线身份、周期、间隙、曲率、实际弧长和铜体积。400、800、1600 轴向采样与 60 项几何检查支持这些几何量的数值一致性，但旧包装器退出码和未保存的内核实际退出保持原状态。24 股、0.100 mm 铜径样张约有 0.1885 mm² 名义铜面积；它说明几何构造，不等于满足至少 5 mm² 的最终束。三维材质、光照与可旋转视角只用于读图。


当前 DTM/iDTM 演示采用加性离散块 \href{report_zh-CN_single.pdf#E06}{(6)}。同一组带权阻抗块只改变排列顺序，不改变总和；因此该实现中出现的均流不能证明不同真实换位路径等效，或证明跨层结构优于传统多级。与文献 [1, 2] 比较时，应核查沿程换位、层级非交换关系、局部场假设和端接，而不能仅复用方法名称。几何公平、模型充分性和性能改善是分别需要证据的三个问题。


\subsection*{身份坐标与分段}
\hypertarget{05.dtm_derivation.heading}{}
\pdfbookmark[1]{身份坐标与分段}{bookmark-05.dtm_derivation.heading}
设 Ps 将固定的物理股线身份映射到第 s 个截面的空间位置。截面电流变为 Ps i，局部电压在回到身份坐标后为 Ps\(^{\mathrm T}\)ZsPs i。串联累加各段 Δzs 即得到总阻抗。因而位置置换与电流身份须分开保存；分段检查 K、2K、4K 的对象是同一条轨迹和同一组物理股线。周期、间隙和弧长检查分别约束几何可实现性与铜用量。


\subsection*{几何采样与等铜检查}
\hypertarget{05.report_geometry.heading}{}
\pdfbookmark[1]{几何采样与等铜检查}{bookmark-05.report_geometry.heading}
轨迹恢复保留物理股线 ID，不按截面排序重新命名。周期检查应分别看位置闭合和身份置换；间隙与曲率检查使用真实单位，弧长由路径采样积分并与加倍采样比较。铜体积与直流电阻以实际股长进入计算，不能把轴向长度直接用于所有螺旋支路。成稿展示的局部截取和绘图抽样在新运行号下记录，原完整中心线继续保留。


等铜对照应注明圆形名义面积、离散实际面积和绝缘占据的外包络。相同名义面积下，较大的弧长会增加单位轴向长度铜体积和电阻路径长度；几何复杂度自身不产生电磁改善证据。640 股候选中心线与 7/24 股构造演示承担不同尺度任务，公开材料应保留这个区别。


\hypertarget{E06}{}\begin{equation}Z_{\mathrm{tot}}=\sum_{s=1}^{K}\Delta z_s P_s^T Z_s P_s\tag{6}\label{eq:E06}\end{equation}
\begin{figure}[!htbp]
\centering
\setcounter{figure}{8}
\refstepcounter{figure}\label{fig:geometry_comparison}
\hypertarget{fig-geometry_comparison}{}
\includegraphics[width=182mm]{../../../composites/zh-CN/geometry_comparison.pdf}
\par\vspace{4pt}{\small\centering 图 9. 股线中心线的三维结构比较。\par}
\vspace{3pt}{\small\raggedright 面板依次展示直束、单级、传统多级局部和跨层结构。 24 股结构使用 0.100 mm 铜径，传统多级面板展示 3 mm 局部。蓝色标记保持选定股线的物理身份。\par}
\end{figure}
\FloatBarrier
\section*{06 有限长度纵向 PEEC}
\hypertarget{chapter-06}{}
\pdfbookmark[0]{06 有限长度纵向 PEEC}{bookmark-chapter-06}
\subsection*{有限长度纵向 PEEC：理论与离散}
\hypertarget{06.peec_method.heading}{}
\pdfbookmark[1]{有限长度纵向 PEEC：理论与离散}{bookmark-06.peec_method.heading}
纵向模型以截面体积细丝承载轴向、沿长度均匀的复电流。欧姆项为各细丝的有限段电阻，部分电感由有限体积对的 1/|r−r′| 核积分得到 \href{report_zh-CN_single.pdf#E08}{(7)}，公共端口和总电流仍按 \href{report_zh-CN_single.pdf#E04}{(4)} 求解。这里的 1264 与 5024 是体积细丝自由度，不是相同数量的绝缘 Litz 股。场内径向电流集中可以由截面电流分配表现，但基函数不含完整横向涡流环路或真实端部重分布。


互项积分阶数与自项精度作为独立选项归档。1264 自由度固定截面首先使用 0.16、0.32、0.64 m 长度；积分检查固定长度，在互项 q=4、6、8 和自项 10/8\(\to\)12/10 之间分别改变一个因素。5024 代表模型复核末两档长度和积分。由于 m40/m80 的实际铜面积略有不同，跨 m 对照按每档自身的等面积 Bessel 参照解释，不能把面积变化隐藏为纯离散误差。


部分电感中的外部磁场与长直圆线的内部 Bessel 阻抗不同，因此这里主要以 Rac 作解析比较，并保留长度和几何面积。长度变化限制了当前有限核计算对所选长度的敏感性，却没有增加新的横向基函数。与 FEM–PEEC 文献 [4] 比较时，本研究的纵向体积细丝、激励、回流与反馈方式应分别注明；某个耦合实现的假设不能推广为所有 PEEC 方法的共同缺点。


\subsection*{基函数与积分}
\hypertarget{06.peec_derivation.heading}{}
\pdfbookmark[1]{基函数与积分}{bookmark-06.peec_derivation.heading}
每条体积细丝内以常量轴向电流密度 iα/Aα 表示电流。把矢势的体积积分投影回测试细丝并沿长度积分，就得到对称部分电感矩阵。自项核可积但在重合点具有奇性，需采用其专用积分精度；互项使用独立阶数。将有限段欧姆电阻与 jωL 相加形成 Z，再由共同端增广系统得到各细丝复电流与端电压。


\subsection*{纵向 PEEC 敏感性与结果}
\hypertarget{06.peec_results.heading}{}
\pdfbookmark[1]{纵向 PEEC 敏感性与结果}{bookmark-06.peec_results.heading}
1264 自由度的 0.16、0.32、0.64 m 三档 Rac 分别为 19.899215、19.958642、19.988477 mΩ/轴向m。两个长度增量的变化均低于 1\%；5024 自由度的 0.32\(\to\)0.64 m 变化也低于门槛。互项 q4/6/8、自项 10/8\(\to\)12/10 以及当前实际截面的粗细离散比较均达到专项规定的有限检查目标。试验逐次改变一个因素；未运行的 1.28 m 档保留为未运行，不能填入外推值或零误差。


末档 5024 自由度、Lz=0.64 m、互项 q4、自项 12/10 的 Rac=20.012738128 mΩ/轴向m、Rac/Rdc=3.644719868，损耗为 8.005095251 W/轴向m；与本档等面积 Bessel 参照的差为 0.044582\%。截面颜色图按保存的细丝电流和方形单元恢复。该细丝电流偏差量描述离散导体内部的电流集中，不是 Litz 物理绝缘股线的均流指标。


长度敏感性与解析参照分别回答不同问题。前者检验有限长度核在所选区间内的变化，后者核验电阻是否与可比较的长直内场尺度一致。二者接近并不消除边界或基函数缺口。特别是末档的很小积分变化只能支持当前积分实现，对真实端区、曲线股线的横向电流以及端接再分布没有新增证据。论文因而把这一结果称为纵向模型内的有限验收。


\hypertarget{E08}{}
\begin{equation}
\begin{gathered}
R_{\alpha\beta}=\delta_{\alpha\beta}\frac{L_z}{\sigma A_\alpha} \\
L_{\alpha\beta}=\frac{\mu}{4\pi A_\alpha A_\beta}\int_{V_\alpha}\int_{V_\beta}\frac{\hat{\mathbf t}_\alpha\cdot\hat{\mathbf t}_\beta}{|\mathbf r-\mathbf r'|}\,\mathrm dV\,\mathrm dV'
\end{gathered}
\tag{7}\label{eq:E08}
\end{equation}
\begin{figure}[!htbp]
\centering
\setcounter{figure}{9}
\refstepcounter{figure}\label{fig:peec_cross_section}
\hypertarget{fig-peec_cross_section}{}
\includegraphics[width=88.9mm]{../../../panels/zh-CN/peec_cross_section.pdf}
\par\vspace{4pt}{\small\centering 图 10. 纵向 PEEC 截面电流。\par}
\vspace{3pt}{\small\raggedright 方形细丝单元的复电流恢复 RMS 幅值。 代表模型含 5024 条纵向体积细丝，轴向长度 0.64 m。坐标和材料面积采用该档实际方形离散定义。\par}
\end{figure}
\begin{figure}[!htbp]
\centering
\setcounter{figure}{10}
\refstepcounter{figure}\label{fig:peec_length}
\hypertarget{fig-peec_length}{}
\includegraphics[width=88.9mm]{../../../panels/zh-CN/peec_length.pdf}
\par\vspace{4pt}{\small\centering 图 11. 纵向 PEEC 长度敏感性。\par}
\vspace{3pt}{\small\raggedright 同一截面按保存的轴向长度比较单位长度电阻。 1264 自由度模型采用 0.16、0.32、0.64 m，5024 自由度复核末两档。积分参数在每个长度序列内固定。\par}
\end{figure}
\begin{figure}[!htbp]
\centering
\setcounter{figure}{11}
\refstepcounter{figure}\label{fig:peec_sensitivity}
\hypertarget{fig-peec_sensitivity}{}
\includegraphics[width=88.9mm]{../../../panels/zh-CN/peec_sensitivity.pdf}
\par\vspace{4pt}{\small\centering 图 12. 纵向 PEEC 积分敏感性。\par}
\vspace{3pt}{\small\raggedright 互项阶数与自项精度各自进行单因素对照。 互项使用 q4、q6、q8，自项使用 10/8 与 12/10。比较保留同一截面及轴向长度。\par}
\end{figure}
\FloatBarrier
\section*{07 RT0 三维外场算例}
\hypertarget{chapter-07}{}
\pdfbookmark[0]{07 RT0 三维外场算例}{bookmark-chapter-07}
RT0 算例固定为 1×1×2 mm 均质铜长方体、200 kHz、Bz=1 mT RMS 外场和零净输运。三分量电流通过面通量自由度表示，连续性、公共端条件与外场功率分别检查。统一 q14 的 8192、16000、27648 单元损耗为 0.927463、0.933539、0.936927 mW；末两次变化为 0.655090\% 和 0.362916\%，均低于 1\%。末档 q18 对照的损耗相对变化为 1.54941×10⁻⁸，复面通量变化为 1.39291×10⁻⁸，低于 0.1\% 积分门槛。


场图由保存的面通量恢复。中心 J 采用相对两面的通量平均除以单元面积；单元平均损耗密度采用完整 RT0 局部重叠积分，不能用中心 |J|²/σ 替代。三维箭头仅显示固定相位零时刻的实部，按统一比例缩放；色标使用复 RMS 范数。切面的实际 z 位置、稀疏矢量抽样和整体积分分别标注，避免把抽样图看成完整场积分或全场最大值。


各档归档的 GMRES 迭代为 35 次，真残差满足 10⁻¹¹ 目标。原 q18 摘要中 null 指标及失败状态仍原样保留，新 CSV 后处理验收另列。旧致密算法的内存拒绝与新版矩阵自由成功求解是不同实现的证据，不能合并为一个状态。当前结论限于固定外场长方体，尚未覆盖曲线 Litz 掩膜、独立边界通量验证或实际输运端接。


\begin{figure}[!htbp]
\centering
\setcounter{figure}{12}
\refstepcounter{figure}\label{fig:rt0_sections}
\hypertarget{fig-rt0_sections}{}
\includegraphics[width=182mm]{../../../composites/zh-CN/rt0_sections.pdf}
\par\vspace{4pt}{\small\centering 图 13. RT0 外场电流与损耗切面。\par}
\vspace{3pt}{\small\raggedright 铜长方体为 1×1×2 mm，200 kHz、1 mT RMS 外场、零净输运。 (a) 单元平均损耗密度；(b) 复 RMS 电流范数。 电流由相对面通量恢复，单元平均损耗采用 RT0 重叠积分。切面位置和色标对应归档的重构定义。\par}
\end{figure}
\begin{figure}[!htbp]
\centering
\setcounter{figure}{13}
\refstepcounter{figure}\label{fig:rt0_vectors}
\hypertarget{fig-rt0_vectors}{}
\includegraphics[width=88.9mm]{../../../panels/zh-CN/rt0_vectors.pdf}
\par\vspace{4pt}{\small\centering 图 14. RT0 重构电流三维矢量。\par}
\vspace{3pt}{\small\raggedright 铜长方体为 1×1×2 mm，200 kHz、1 mT RMS 外场、零净输运。 箭头为固定相位的实部，色标为复 RMS 范数。 矢量采用固定比例与稀疏单元采样。单元与面通量身份在保存的重构数据中对应。\par}
\end{figure}
\begin{figure}[!htbp]
\centering
\setcounter{figure}{14}
\refstepcounter{figure}\label{fig:rt0_spatial}
\hypertarget{fig-rt0_spatial}{}
\includegraphics[width=88.9mm]{../../../panels/zh-CN/rt0_spatial.pdf}
\par\vspace{4pt}{\small\centering 图 15. RT0 空间与积分对照。\par}
\vspace{3pt}{\small\raggedright 铜长方体为 1×1×2 mm，200 kHz、1 mT RMS 外场、零净输运。 q14 空间序列与末档 q18 对照使用相同外场工况。 空间序列为 8192、16000、27648 单元。积分对照在最终成功网格上进行。\par}
\end{figure}
\FloatBarrier
\section*{08 三维 A–φ 场与诊断}
\hypertarget{chapter-08}{}
\pdfbookmark[0]{08 三维 A–φ 场与诊断}{bookmark-chapter-08}
A–φ 电流应使用实际保存的 A 和 φ 梯度恢复 \href{report_zh-CN_single.pdf#E07}{(8)}，不能把未通过的 φ 替换成解析值。三档圆柱分别为 1683、3672、9108 节点。最细档的径向 J 校准达到其门槛，但 φ 与规范残差在原三档中仍未通过；较粗档 J 本身也未通过。制造解核验检验了指定方程与实现关系，却没有使这些圆柱物理解自动通过。


诊断图使用 128 个具有明确节点身份的保存场及导数样本，分别显示 A、φ、J 和实际 φ 梯度。它们不是从美化表面反推的结果，也不是全场极值。边界、缩放、导数和后处理诊断应区分实现问题与离散误差；达到有限修订上限后，未通过项以触发条件、最小复现和继续研究所需条件封存。直接分解、OOC 与迭代失败路线也保持各自原因。


\hypertarget{E07}{}\begin{equation}\mathbf J=-\sigma\left(j\omega\mathbf A+\nabla\phi\right)\tag{8}\label{eq:E07}\end{equation}
\begin{figure}[!htbp]
\centering
\setcounter{figure}{15}
\refstepcounter{figure}\label{fig:aphi_preserved}
\hypertarget{fig-aphi_preserved}{}
\includegraphics[width=182mm]{../../../composites/zh-CN/aphi_preserved.pdf}
\par\vspace{4pt}{\small\centering 图 16. A–φ 保存场样本。\par}
\vspace{3pt}{\small\raggedright 面板分别给出 A、φ、J 和实际 φ 梯度。 128 个指定节点保留其样本身份。色标使用实际保存的复数场与导数幅值。\par}
\end{figure}
\begin{figure}[!htbp]
\centering
\setcounter{figure}{16}
\refstepcounter{figure}\label{fig:aphi_gates}
\hypertarget{fig-aphi_gates}{}
\includegraphics[width=88.9mm]{../../../panels/zh-CN/aphi_gates.pdf}
\par\vspace{4pt}{\small\centering 图 17. A–φ 圆柱校准指标。\par}
\vspace{3pt}{\small\raggedright 三档网格与各指标门槛分别展示。 圆柱三档含 1683、3672、9108 节点。误差以各自验收门槛归一化，水平线标出单位比值。\par}
\end{figure}
\FloatBarrier
\section*{09 Ansys 算例与接口验证}
\hypertarget{chapter-09}{}
\pdfbookmark[0]{09 Ansys 算例与接口验证}{bookmark-chapter-09}
七股二维历史有效端口短测使用 704 单元，包含两个电压基解和实际 ±20 A 场解；公共端电压、逐股复电流与独立功率支持该单档算例。它不是三档网格收敛。新固定网格与精度尝试未形成要求的完整验收；原生三维 J 表达式探针也未得到可接受物理解，因此不为其生成场结果图或填造数值。


归档边界采用已核查 v252 Student 的二维 2000 三角形、三维 64000 体单元/8000 面单元及四核限制。接口错误、网格或许可限制和数值误差有不同含义，应分别记录。实际端接、曲线 STL 和足铜大束继续需要相应许可、网格能力及独立验证；有限短测不能用于推荐真实换位性能。


\begin{figure}[!htbp]
\centering
\setcounter{figure}{17}
\refstepcounter{figure}\label{fig:ansys_seven_current}
\hypertarget{fig-ansys_seven_current}{}
\includegraphics[width=88.9mm]{../../../panels/zh-CN/ansys_seven_current.pdf}
\par\vspace{4pt}{\small\centering 图 18. Ansys 七股二维复电流。\par}
\vspace{3pt}{\small\raggedright 704 单元有效端口短测给出逐股幅值与相位。 电流采用实际 ±20 A 场解的端口记录。幅值与相位对应同一七股编号。\par}
\end{figure}
\begin{figure}[!htbp]
\centering
\setcounter{figure}{18}
\refstepcounter{figure}\label{fig:ansys_limits}
\hypertarget{fig-ansys_limits}{}
\includegraphics[width=88.9mm]{../../../panels/zh-CN/ansys_limits.pdf}
\par\vspace{4pt}{\small\centering 图 19. Ansys Student 能力与接口诊断。\par}
\vspace{3pt}{\small\raggedright Ansys Student v252 的二维三角形、三维体单元与面单元上限分别为 2000、64000 和 8000，计算使用最多 4 核。并列分组表示各项许可限制。电流密度表达式 J 的接口探针与许可规模分别记录，实际求解状态见三维诊断。\par}
\end{figure}
\FloatBarrier
\section*{10 计算实现、CPU/GPU 对照与资源表现}
\hypertarget{chapter-10}{}
\pdfbookmark[0]{10 计算实现、CPU/GPU 对照与资源表现}{bookmark-chapter-10}
\subsection*{CPU/GPU 等价性与资源边界}
\hypertarget{10.compute.heading}{}
\pdfbookmark[1]{CPU/GPU 等价性与资源边界}{bookmark-10.compute.heading}
CPU/GPU 对照固定同一纵向模型和离散设置。5024 自由度的 L 矩阵相对差为 7.48149×10⁻¹⁶，复电流相对差为 7.78082×10⁻¹⁴，低于 10⁻¹⁰ 门槛；1264 对照亦通过。这支持已实现 CPU/GPU 算法在所比较矩阵与端口解上的等价。它不赋予模型额外物理内容，也不能用不同规模或不同积分设置的用时声称加速比。


9856 自由度历史结果只恢复了部分电流等数据，缺少完整矩阵、公共电压或全场的项目不升级为完整验证。20108 及更大规模保持已有资源边界证据和退出记录。资源曲线只采用可归属的进程树样本；短时峰值可能遗漏，混入其他 PID 的记录排除。计算能力边界与物理失败分开解释：缺少资源不是通过，运行完成也不是物理验收通过。


\subsection*{构建与复现入口}
\hypertarget{10.workflow.heading}{}
\pdfbookmark[1]{构建与复现入口}{bookmark-10.workflow.heading}
复现分为研究计算与成稿构建两个入口。研究入口读取原始输入和求解配置，执行相应验证；成稿入口读取冻结数据、已验收单图与 Markdown 正文，生成章节、合并稿和多种格式。每一阶段保存输入哈希和输出哈希，只有依赖与产物一致时才复用已完成阶段。独立单图保留矢量曲线及其字体，组图清单保存面板身份与布局参数。


\begin{figure}[!htbp]
\centering
\setcounter{figure}{19}
\refstepcounter{figure}\label{fig:gpu_equivalence}
\hypertarget{fig-gpu_equivalence}{}
\includegraphics[width=88.9mm]{../../../panels/zh-CN/gpu_equivalence.pdf}
\par\vspace{4pt}{\small\centering 图 20. CPU/GPU 同模型数值对照。\par}
\vspace{3pt}{\small\raggedright 矩阵与复电流相对差使用同一离散和积分设置。 相对差针对 1264 和 5024 自由度分别计算。水平参照为 10⁻¹⁰ 算法对照门槛。\par}
\end{figure}
\begin{figure}[!htbp]
\centering
\setcounter{figure}{20}
\refstepcounter{figure}\label{fig:resource_memory}
\hypertarget{fig-resource_memory}{}
\includegraphics[width=88.9mm]{../../../panels/zh-CN/resource_memory.pdf}
\par\vspace{4pt}{\small\centering 图 21. 可归属进程树资源曲线。\par}
\vspace{3pt}{\small\raggedright 曲线使用归档的时间及内存样本。 时间以对应运行的起点归一化。内存为可归属进程树的采样工作集。\par}
\end{figure}
\FloatBarrier
\section*{11 综合讨论}
\hypertarget{chapter-11}{}
\pdfbookmark[0]{11 综合讨论}{bookmark-chapter-11}
\subsection*{讨论：模型不确定性与文献关系}
\hypertarget{11.discussion.heading}{}
\pdfbookmark[1]{讨论：模型不确定性与文献关系}{bookmark-11.discussion.heading}
结果表明，均流与低损耗不是只由单股尺寸决定的同一问题。共端模型允许电流的幅值和相位自洽变化，而当前几何展示只提供轨迹与弧长代价。真实性能比较须把这种共端约束带入具有充分沿程耦合的电磁模型，并采用相同的实际铜截面、回流和长度定义。候选改善还应大于相关数值不确定性；未解决的模型误差不能被很小的求解残差抵消。


与 [1, 2] 的关系主要在换位表示及其适用范围，而不是本文已经复现全部实验。与 [3] 的关系在可解释的降阶电路，但其特定绕组的实验误差不能移作本项目精度。[5] 的公共电压与总电流约束提供方法联系，其频率、尺寸及准二维假设不直接验证本项目 200 kHz 大束。均匀化文献 [6, 7] 则提示宏观场与微观局部损耗需要一致分配，避免重复计数；平均化并不会自动证明股间电流均匀。


本研究没有热模型或实验测量来检验绕制误差、端接接触、绝缘厚度和温升。RT0 外场收敛与纵向 PEEC 电阻接近解析分别提供局部证据；未通过三维 A–φ 与 Ansys 原生接口留下实际研究边界。后续优先任务应是具有真实端部、回流与横向自由度的可承受小股数三维交叉验证，再扩展到足铜多级强基线和候选结构，最后进行受控实测。


\subsection*{后续验证顺序}
\hypertarget{11.recommendations.heading}{}
\pdfbookmark[1]{后续验证顺序}{bookmark-11.recommendations.heading}
后续研究宜按模型缺口安排验证顺序：先在可承受的三维校准体上建立边界、规范与端场的一致性，再引入小股数真实端接和横向自由度，最后比较足铜传统多级与跨层结构。实验应在相同频率、温度、回流和轴向长度下测量阻抗，并记录端接电阻与几何公差。候选改进的幅度应与数值及测量不确定性共同比较。


\FloatBarrier
\section*{12 结论}
\hypertarget{chapter-12}{}
\pdfbookmark[0]{12 结论}{bookmark-chapter-12}
第一，分层验证框架把共同物理内容与各层缺口分别记录，使解析、二维、纵向及三维外场结果能够在可比较范围内使用。第二，共端增广系统给出了指定线性模型内的复电流均流判据，避免用预设均流替代验证。第三，有限长度纵向 PEEC 的已运行长度与积分敏感性达到专项门槛，其末档电阻与本档等面积解析参照相符；该结论继续受纵向基函数范围限制。


640 股二维场、固定外场长方体 RT0 和 CPU/GPU 同模型对照形成可重绘的冻结证据。实际换位束的三维电磁排名、热模型和实验仍未收口；失败路线保留触发条件、原始结果和后续所需能力。


\FloatBarrier
\section*{A 数学推导}
\hypertarget{chapter-A}{}
\pdfbookmark[0]{A 数学推导}{bookmark-chapter-A}
\subsection*{解析归一化、场路与单元积分推导}
\hypertarget{A.report_derivation.heading}{}
\pdfbookmark[1]{解析归一化、场路与单元积分推导}{bookmark-A.report_derivation.heading}
圆线解析归一化先要求 2π\(\int\)\(_0\)\(^a\) r Jz(r) dr=It，再用 d[r I\(_1\)(κr)]/dr=κr I\(_0\)(κr) 得到分母。边界轴向电场 Jz(a)/σ 除以 It 给出内部单位长度阻抗。独立焦耳积分与 |It|² Re(zint) 对照采用 RMS 约定；高精度积分误差仅检验所写解析公式。有限端部、回流外部电感、热变化和绝缘结构并不包含在该解析精度中。


线性三角形场恢复使用实际物理面积 aT。令节点 eₖ=v−uₖ，则单元复电流为 σaT\(\sum\)eₖ/3，焦耳功率为 σaT[\(\sum\)|eₖ|²+|\(\sum\)eₖ|²]/12。后者来自对角系数 2、非对角系数 1 的一致质量矩阵。原坐标以 0.001 m 缩放，因此面积乘缩放平方；显示到 mm 不影响积分。材料标记区分每股、回流和空气，回流使用自身 v。


共端判据的必要性和充分性依赖同一个线性系统及解唯一性。必要性来自均流代入端压方程；充分性来自被 Q 消去的向量可以表示为 V1，并与总电流约束共同满足增广系统。导体接触、额外端接阻抗、非线性材料或遗漏的感应项会改变系统。因此在新模型里复用判据时，先重建 Z、e 与约束，而不是只复用结果。


\subsection*{判据推导}
\hypertarget{A.sharing_derivation.heading}{}
\pdfbookmark[1]{判据推导}{bookmark-A.sharing_derivation.heading}
均流判据的充分性来自端口解的唯一性。当投影条件成立时，向量 (It/N)Z1+e 的各分量相同，因此存在共同电压使等流向量满足全部端口方程及总电流约束。必要性则由已存在的等流解直接代入得到。数值计算以归一化复电流偏差评估均流程度，同时报告幅值和相位，以保持复数求和的信息。


\FloatBarrier
\section*{B 离散与后处理}
\hypertarget{chapter-B}{}
\pdfbookmark[0]{B 离散与后处理}{bookmark-chapter-B}
\subsection*{独立单元积分}
\hypertarget{B.triangle_derivation.heading}{}
\pdfbookmark[1]{独立单元积分}{bookmark-B.triangle_derivation.heading}
线性三角形上 J=\(\sum\)\(_a\)NaJa。复电流积分使用每个形函数积分 A/3；焦耳积分使用质量矩阵与焦耳积分见 \href{report_zh-CN_single.pdf#E09}{(9)}。材料标记决定本单元所属股线、回流或空气。坐标先转换为米再计算面积，所得总功率与原归档标量进行独立一致性检查。


\subsection*{RT0 面通量与轻量场恢复}
\hypertarget{B.report_rt0.heading}{}
\pdfbookmark[1]{RT0 面通量与轻量场恢复}{bookmark-B.report_rt0.heading}
原 RT0 面通量按方向及整数网格身份保存。单元每轴取相对两面的 q−、q+，中心分量为 (q−+q+)/(2h²)，轴向局部电阻功率为 [|q−|²+|q+|²+Re(q−*q+)]/(3σh)。三轴相加、全单元求和恢复总损耗。被绝缘侧边省略的面在原离散定义中具有结构零通量；这种零值不等于另做了独立边界通量场解。


成稿后处理独立恢复全部 27648 单元中心复 J，与原 cell\_fields.csv 对照；同时用完整重叠积分核验原损耗。切面取实际最接近中心的离散层并标注 z，不假设所有采样正好位于 1 mm。W/cm³ 色标由 SI 损耗密度乘 10⁻⁶ 得到。三维箭头的抽样身份与公共长度缩放可在图件请求和映射数据中检查，不用额外求解器。


\hypertarget{E09}{}
\begin{equation}
\begin{gathered}
M_T=\frac{A_T}{12}\begin{bmatrix}2&1&1\\1&2&1\\1&1&2\end{bmatrix} \\
P_T=\mathbf J_T^H M_T\mathbf J_T/\sigma
\end{gathered}
\tag{9}\label{eq:E09}
\end{equation}
\FloatBarrier
\section*{C 验收、诊断与复现}
\hypertarget{chapter-C}{}
\pdfbookmark[0]{C 验收、诊断与复现}{bookmark-chapter-C}
每条路线保留模型、参数、源码和运行身份，原始 WXF、检查点、Ansys 工程、日志及失败文件不被转换结果覆盖。JSON/CSV 提供复数实虚部、单位和逐股身份；原字段的单位标注若存在歧义，由源模型定义和附加说明解释，不改写原文件。逐图映射记录来源数据、表达式、独立 PDF/SVG/PNG 哈希及组图布局。有效场积分与完整空间场导出分开记录，部分恢复不填造成完整场。


失败最小复现应包含触发输入、实际错误、尝试过的针对性修正、停止条件及下一步所需资源或许可。已知相同代码和输入的接口错误不通过重复求解变成“更多证据”。已有源文件和日志通过新清单指向，既不删除历史异常退出，也不把后处理成功回填为原求解成功。这样才能把研究限制交给其他研究者继续验证。


成稿运行限制 CPU 线程、任务内存和系统空闲内存，只读解析缓存并后处理。Notebook 保存静态正文、公式与图形，原生 Graphics3D 允许前端旋转，但打开文稿不运行 FEM、PEEC 或 GPU 求解。缺少附件时报告具体文件。带空格恢复副本、哈希一致性和轻量读取是移交验收的一部分；正文排版与图件导出通过不改变物理验收状态。


\begin{table}[!htbp]
\centering{\footnotesize q1\_fem2\_convergence.csv}\par
\begingroup\footnotesize\begin{tabular}{p{40.45mm}|p{40.45mm}|p{40.45mm}|p{40.45mm}}
\toprule
nodes & h\_\allowbreak over\_\allowbreak a & Rac\_\allowbreak ohm\_\allowbreak per\_\allowbreak m & Rac\_\allowbreak relative\_\allowbreak error\\
353 & 0.24 & 0.02015722 & 0.007045247\\
1389 & 0.12 & 0.02002081 & 0.0002302695\\
5573 & 0.06 & 0.02001617 & 1.79233e-06\\
\bottomrule\end{tabular}\endgroup
\end{table}
\begin{table}[!htbp]
\centering{\footnotesize peec\_length.csv}\par
\begingroup\footnotesize\begin{tabular}{p{40.45mm}|p{40.45mm}|p{40.45mm}|p{40.45mm}}
\toprule
filaments & length\_\allowbreak m & Rac\_\allowbreak ohm\_\allowbreak per\_\allowbreak axial\_\allowbreak m & Rac\_\allowbreak over\_\allowbreak Rdc\\
1264 & 0.16 & 0.01989922 & 3.647128\\
1264 & 0.32 & 0.01995864 & 3.65802\\
1264 & 0.64 & 0.01998848 & 3.663488\\
5024 & 0.32 & 0.01998302 & 3.639307\\
5024 & 0.64 & 0.02001274 & 3.64472\\
\bottomrule\end{tabular}\endgroup
\end{table}
\begin{table}[!htbp]
\centering{\footnotesize rt0\_cases.csv}\par
\begingroup\footnotesize\begin{tabular}{p{84.9mm}|p{84.9mm}}
\toprule
cells & recomputed\_\allowbreak loss\_\allowbreak W\\
8192 & 0.0009274631\\
16000 & 0.0009335389\\
27648 & 0.0009369268\\
27648 & 0.0009369268\\
\bottomrule\end{tabular}\endgroup
\end{table}
\begin{table}[!htbp]
\centering{\footnotesize gpu\_equivalence.csv}\par
\begingroup\footnotesize\begin{tabular}{p{55.26666666666667mm}|p{55.26666666666667mm}|p{55.26666666666667mm}}
\toprule
filaments & L\_\allowbreak matrix\_\allowbreak relative\_\allowbreak error & current\_\allowbreak relative\_\allowbreak error\\
1264 & 7.075722e-16 & 3.281741e-14\\
5024 & 7.481492e-16 & 7.780817e-14\\
\bottomrule\end{tabular}\endgroup
\end{table}
\FloatBarrier
\section*{D 数据来源与环境交接}
\hypertarget{chapter-D}{}
\pdfbookmark[0]{D 数据来源与环境交接}{bookmark-chapter-D}
复现分为研究计算与成稿构建两个入口。研究入口读取原始输入和求解配置，执行相应验证；成稿入口读取冻结数据、已验收单图与 Markdown 正文，生成章节、合并稿和多种格式。每一阶段保存输入哈希和输出哈希，只有依赖与产物一致时才复用已完成阶段。独立单图保留矢量曲线及其字体，组图清单保存面板身份与布局参数。


来源清单 / Source manifest: source\_provenance.json; quantities.json; asset\_sources.json.


\begin{table}[!htbp]
\centering{\footnotesize 冻结来源与完整性 / Frozen sources and integrity}\par
\begingroup\footnotesize\begin{tabular}{p{84.9mm}|p{84.9mm}}
\toprule
Source & SHA256\\
docs/总控/当前运行索引.json & \seqsplit{3324976d4af8f58e60820a88ae978246443f78fbdd0cfa1050e86809110ae60d}\\
outputs/prewrite\_\allowbreak closure\_\allowbreak state.json & \seqsplit{80a66afa79c321a89f730b102af0e70125f6e1655372717d6b59f470c6520e71}\\
\bottomrule\end{tabular}\endgroup
\end{table}
\clearpage
\begin{thebibliography}{99}
\bibitem[1]{P01} Silvano Cruciani, Tommaso Campi, Francesca Maradei, Mauro Feliziani. Numerical Modeling of Litz Wires Based on Discrete Transpositions of Strands and 2-D Finite Element Analysis. IEEE TPEL 38(5):6710–6719, 2023. https://doi.org/10.1109/TPEL.2023.3240338
\bibitem[2]{N01} Xinsheng Zhang, Fei Xiao, Ruitian Wang, Shunjun Chen, Xuexin Fan, Xin Li. AC Impedance Modeling of Composite Litz Wire With Nonideal Transposition Based on Improved Discrete Transposition Model. IEEE TPEL 41(7):11003–11019, 2026. https://doi.org/10.1109/TPEL.2026.3661080
\bibitem[3]{P02} Yueyin Wang, Wu Chen, Zhan Shen. Modeling of Litz Wire Winding AC Resistance Based on Improved Segmented Equivalent Circuit. IEEE TPEL 39(1):1135–1149, 2024. https://doi.org/10.1109/TPEL.2023.3321327
\bibitem[4]{P03} Andreas Roßkopf, Eberhard Bär, Christopher Joffe, Clemens Bonse. Calculation of Power Losses in Litz Wire Systems by Coupling FEM and PEEC Method. IEEE TPEL 31(9):6442–6449, 2016. https://doi.org/10.1109/TPEL.2015.2499793
\bibitem[5]{P04} Tamás Bakondi, Sándor Bilicz, Szabolcs Gyimóthy. Finite element modeling of eddy currents in a multiconductor system with quasi-2D symmetry. COMPEL 44(5):730–742, 2025. https://doi.org/10.1108/COMPEL-01-2025-0031
\bibitem[6]{P05} David C. Meeker. An Improved Continuum Skin and Proximity Effect Model for Hexagonally Packed Wires. JCAM 236(18):4635–4644, 2012. https://doi.org/10.1016/j.cam.2012.04.009
\bibitem[7]{P09} Korawich Niyomsatian, Johan Gyselinck, Ruth V. Sabariego. Closed-form complex permeability expression for proximity-effect homogenisation of litz-wire windings. IET SMT 14(3):287–291, 2020. https://doi.org/10.1049/iet-smt.2019.0377
\end{thebibliography}
\end{document}
